\documentclass[11pt]{article}

\usepackage[margin=0.65in]{geometry}
\usepackage{amsmath}
\usepackage{booktabs}
\usepackage{caption}
\usepackage{float}
\usepackage{graphicx}
\usepackage{hyperref}
\usepackage{longtable}
\usepackage{xcolor}

\setlength{\parindent}{0pt}
\setlength{\parskip}{0.35em}
\captionsetup{font=small}

\title{Executive Summary: Physics Infrastructure for the Asymmetric Microfluidic Loop}
\author{Physics-Embedded Transformers Project}
\date{July 22, 2026}

\begin{document}

\maketitle

\begin{abstract}
This document summarizes the physics-only portion of the project: device
geometry, region labeling, droplet occupancy, branch hydraulics, full-device
CFD, flow-split calibration, and interpolation of the computed continuous-phase
velocity field onto measured droplet trajectories. It intentionally avoids
machine-learning architecture and training details. The present physics
pipeline produces a geometrically clean full-device computational domain, a
validated finite-element Stokes solution on the frozen 24~\(\mu\mathrm{m}\)
production mesh, a calibrated library of flow splits from 0.40 to 0.64 left
fraction, and CFD-enriched droplet tracking data for video 2.
\end{abstract}

\section{Physical System}

The device is an asymmetric loop microfluidic channel. A single inlet channel
feeds an upstream junction, where the flow separates into a long left branch
and a shorter right branch. The two branches reconnect at a downstream junction
and discharge through a single outlet channel. The channel width and height are
both 100~\(\mu\mathrm{m}\), corresponding to 25 pixels at the calibrated
4~\(\mu\mathrm{m/px}\) image scale.

The measured branch lengths are:
\[
L_{\mathrm{left}} = 1791~\mu\mathrm{m}, \qquad
L_{\mathrm{right}} = 1491~\mu\mathrm{m}.
\]
The resulting geometric asymmetry is:
\[
\frac{L_{\mathrm{left}}}{L_{\mathrm{right}}} = 1.201207,
\qquad
\alpha_{\mathrm{geom}}
= \frac{L_{\mathrm{left}} - L_{\mathrm{right}}}{L_{\mathrm{right}}}
= 0.201207.
\]

The working physical picture is that droplets occupy finite portions of the
network and alter the effective branch resistance. The continuous phase is
treated as a low-Reynolds-number background flow field, while droplet occupancy
and branch hydraulics provide framewise physical state variables.

\section{Physics Tools}

The physics pipeline is organized by scientific responsibility rather than by
implementation history:

\begin{enumerate}
  \item \textbf{Geometry and region labeling:} reconstruct the calibrated
        device coordinate system, centerlines, branches, junctions, and
        channel-region labels.
  \item \textbf{Droplet occupancy:} compute fractional droplet occupancy in
        physical regions of the device.
  \item \textbf{Branch hydraulics:} estimate framewise branch resistance and
        flow split from normalized branch occupancy.
  \item \textbf{Full-device CFD:} solve the steady continuous-phase Stokes
        problem over the full loop geometry.
  \item \textbf{CFD library and interpolation:} calibrate a library of flow
        splits and sample the corresponding velocity field at arbitrary valid
        points in the device.
  \item \textbf{Physics-enriched tracking:} append local CFD velocity and
        hydraulic state to measured droplet centroids.
\end{enumerate}

These tools are deterministic preprocessing and simulation components. Their
outputs are reusable physical data products.

\section{Geometry and Mesh}

The full device is constructed as a continuous CAD-style fluid region using the
known centerline topology and the fixed 100~\(\mu\mathrm{m}\) channel width.
The domain is represented in physical units, with coordinates stored in
micrometers. The final computational domain has one connected fluid component
and one central island. Dedicated topology checks confirmed no self
intersections, no dangling boundaries, no duplicated interior wall fragments,
and no unintended internal boundaries.

\begin{figure}[H]
  \centering
  \includegraphics[width=0.48\linewidth]{../outputs/physics/full_device_cfd/mesh/full_device_polygon_with_centerlines.png}
  \includegraphics[width=0.48\linewidth]{../outputs/physics/full_device_cfd/mesh/full_device_mesh_quality.png}
  \caption{Left: full-device CAD polygon with centerlines. Right: mesh-quality
  visualization for the frozen production mesh.}
  \label{fig:full-device-geometry}
\end{figure}

The frozen production mesh uses the 24~\(\mu\mathrm{m}\) configuration. Mesh
quality is summarized in Table~\ref{tab:mesh-quality}. The mesh is fine enough
to resolve the junctions and curved walls while avoiding unnecessary refinement
in the long straight sections.

\begin{table}[H]
\centering
\caption{Frozen 24~\(\mu\mathrm{m}\) production mesh quality.}
\label{tab:mesh-quality}
\begin{tabular}{@{}lr@{}}
\toprule
Quantity & Value \\
\midrule
Nodes & 2339 \\
Triangular elements & 3807 \\
Minimum angle & \(25.73^\circ\) \\
Maximum angle & \(122.39^\circ\) \\
Median aspect ratio & 1.456 \\
Maximum aspect ratio & 2.285 \\
Fluid area & \(5.226 \times 10^5~\mu\mathrm{m}^2\) \\
Boundary length & \(1.043 \times 10^4~\mu\mathrm{m}\) \\
\bottomrule
\end{tabular}
\end{table}

\section{Continuous-Phase CFD}

The continuous phase is modeled as a steady, incompressible, creeping flow.
The governing equations are the Stokes equations:
\[
-\nabla p + \mu \nabla^2 \mathbf{u} = \mathbf{0},
\qquad
\nabla \cdot \mathbf{u} = 0,
\]
where \(\mathbf{u}\) is the two-dimensional velocity field, \(p\) is pressure,
and \(\mu\) is the dynamic viscosity of the continuous phase. No-slip conditions
are imposed on solid walls. The inlet carries a fixed incoming flux. Outlet
and branch-flow conditions are formulated so that calibrated left/right flow
fractions can be generated without changing the geometry.

The validated equal-split case corresponds to
\[
f_L = \frac{Q_L}{Q_L + Q_R} \approx 0.50.
\]
For this case, the left and right branch fluxes are equal to numerical
precision, and the velocity and pressure fields are smooth over the full device.
The maximum speed in the production-library cases is approximately
0.0817~\(\mathrm{m/s}\). Pressure ranges over the calibrated cases are about
173--190~Pa, increasing as the imposed branch bias is made stronger.

\begin{figure}[H]
  \centering
  \includegraphics[width=0.42\linewidth]{../outputs/physics/full_device_cfd/alpha0_equal_split_smoke/diagnostics/speed_magnitude.png}
  \includegraphics[width=0.42\linewidth]{../outputs/physics/full_device_cfd/alpha0_equal_split_smoke/diagnostics/pressure.png}
  \caption{Representative equal-split full-device CFD solution: velocity
  magnitude and pressure field.}
  \label{fig:cfd-fields}
\end{figure}

\begin{figure}[H]
  \centering
\includegraphics[width=0.48\linewidth]{../outputs/physics/full_device_cfd/alpha0_equal_split_smoke/diagnostics/full_device_streamlines_very_dense_postprocessed.png}
\includegraphics[width=0.48\linewidth]{../outputs/physics/full_device_cfd/library/requested_vs_achieved_split.png}
\caption{Left: dense streamline post-processing for the equal-split case.
Right: production calibration curve for requested versus achieved left-flow
fraction.}
\label{fig:streamlines-and-library}
\end{figure}

\section{Flow-Split Calibration Library}

The production CFD library spans requested left-flow fractions from 0.40 to
0.64 in increments of 0.02. The achieved split is used as the interpolation
coordinate. Calibration succeeded for all 13 requested targets, with no
unreachable cases. The natural unmodified split is
\[
f_{L,\mathrm{natural}} = 0.458119.
\]

The maximum absolute calibration error over the production library is
\[
3.61 \times 10^{-9},
\]
and the achieved splits are strictly monotonic. The largest adjacent gap is
0.020000004, matching the intended resolution.

\begin{table}[H]
\centering
\caption{Production CFD split-library summary.}
\label{tab:split-library}
\begin{tabular}{@{}lrl@{}}
\toprule
Quantity & Value & Interpretation \\
\midrule
Successful cases & 13 & All requested targets reached \\
Requested range & 0.40--0.64 & Left-flow fraction \\
Natural split & 0.458119 & Unbiased full-device flow split \\
Maximum calibration error & \(3.61\times10^{-9}\) & Achieved vs. requested split \\
Largest adjacent gap & 0.020000004 & Library spacing \\
Monotonic achieved splits & yes & Suitable for interpolation \\
\bottomrule
\end{tabular}
\end{table}

\section{Mesh and Field Verification}

A mesh convergence comparison was performed between the frozen 24~\(\mu\mathrm{m}\)
production mesh and a finer 12~\(\mu\mathrm{m}\) mesh. The goal was not to
replace the production mesh, but to verify that the velocity field, pressure
field, and flux structure are stable at the selected resolution. The largest
differences occur near the inlet and outlet junctions, which is expected because
these are the regions with the strongest geometric curvature and flow
redistribution.
The production mesh was retained as the frozen default after this verification:
it captured the same qualitative field structure as the finer mesh while
remaining inexpensive enough for production library generation.

\section{Physics-Enriched Droplet Data}

The full-device CFD library is sampled at measured droplet centroids using the
framewise hydraulic split as the library coordinate. This attaches local
continuous-phase velocity information to each observed droplet without
extrapolating outside the CFD fluid domain.

For video 2, the enriched tracking table contains 432,460 observed droplet rows.
Of these, 429,929 rows are inside the valid CFD fluid domain:
\[
\mathrm{valid~fraction} = 0.994147.
\]
The remaining invalid rows are confined to the outlet channel and are consistent
with the known downstream video crop/metadata exclusion rather than a CFD
failure.

\begin{table}[H]
\centering
\caption{Regional validity of CFD samples on measured droplet centroids.}
\label{tab:regional-validity}
\begin{tabular}{@{}lrrr@{}}
\toprule
Region & Rows & Invalid rows & Invalid percentage \\
\midrule
Inlet channel & 65,614 & 0 & 0.00\% \\
Inlet junction & 26,553 & 0 & 0.00\% \\
Left branch & 100,083 & 0 & 0.00\% \\
Right branch & 170,729 & 0 & 0.00\% \\
Outlet junction & 8,399 & 0 & 0.00\% \\
Outlet channel & 61,082 & 2,531 & 4.14\% \\
\bottomrule
\end{tabular}
\end{table}

\begin{figure}[H]
  \centering
  \includegraphics[width=0.48\linewidth]{../outputs/physics/video_2/enrichment/diagnostics/sampled_field_overlay.png}
  \includegraphics[width=0.48\linewidth]{../outputs/physics/video_2/enrichment/diagnostics/final_validity/invalid_percentage_by_region.png}
  \caption{Left: sampled CFD vectors over tracked droplet centroids. Right:
  invalid CFD percentage by physical region. Invalid samples are isolated to
  the cropped outlet-channel region.}
  \label{fig:enrichment}
\end{figure}

\section{Current Physics Status}

The physics infrastructure is in a production-ready state for the current
single-device, single-experiment analysis:

\begin{itemize}
  \item the full-device geometry is topologically clean;
  \item the 24~\(\mu\mathrm{m}\) mesh is frozen as the production mesh;
  \item the full-device Stokes solve produces smooth pressure and velocity
        fields;
  \item the calibrated CFD library provides monotonic flow splits from 0.40 to
        0.64;
  \item the library is interpolation-ready using achieved left-flow fraction;
  \item measured droplet tracks are enriched with valid local CFD velocity for
        99.4\% of observed rows.
\end{itemize}

Future physics extensions can build on this foundation by adding explicit
droplet-induced resistance laws, pressure-loss models at the junctions,
time-dependent coupling, or additional experimental devices. These are natural
physics-modeling extensions and do not require revisiting the validated
geometry, mesh, and production-library foundation summarized here.

\end{document}
